\section{\textsc{Dream}: Anchor-Modulated Representation}
\label{sec:dream}

Given the entanglement diagnosis of Sec.~\ref{sec:ceiling}, the
natural intervention is to give each subject a personalized
\emph{reference} that captures their identity-correlated baseline so
that downstream readouts can operate on the residual variation. Prior
attempts at per-subject calibration (subtracting each subject's mean
feature~\cite{tian2024tfic}) destroy signal because the mean over a
subject's clips entangles their engagement distribution with their
identity. We sidestep this by introducing a self-supervised,
engagement-agnostic criterion for selecting \emph{anchor frames}.

\subsection{Anchor selection}
\label{sec:anchor-select}

For each frame $x$ we compute MediaPipe's 52 blendshape activations
$b(x) \in [0,1]^{52}$ and 6-d eye gaze $g(x) \in \mathbb{R}^6$ (left
and right yaw, pitch, openness). We score each frame by a
\emph{neutral score}:
\begin{equation}
\mathrm{neutral}(x) \;=\; \|b(x)\|_1 \;+\; \lambda \cdot \|g(x) - g^{\mathrm{fwd}}\|_2,
\end{equation}
where $g^{\mathrm{fwd}} = (0,0,1,0,0,1)$ is the canonical
forward-gaze reference and $\lambda = 0.5$. Lower score $\Rightarrow$
more neutral face (closed mouth, neutral eyes, eyes facing forward).
The score is engagement-agnostic by construction: no engagement label
enters its computation.

For each subject $s$ we select the $K \in \{1, 3, 5, 10, 20\}$ frames
with lowest neutral score as that subject's \emph{anchor set} $A_s$.
For unseen test subjects we apply the same procedure to that
subject's own \emph{unlabelled} clips---i.e., we use the test split's
blendshapes but never its engagement labels. This is operationally
equivalent to a deployment scenario in which a teacher records a
student briefly before scoring begins.

\paragraph{Anchor-label distribution.} Tab.~\ref{tab:anchor-dist}
reports the engagement-label distribution of the selected anchor
frames at $K = 5$ across protocols. Anchors are dominated by $L2$
and $L3$ (the majority classes), matching the overall test-set
distribution---confirming that neutral-face frames are
\emph{behaviour-balanced}, not biased toward low-engagement clips.
The criterion is genuinely engagement-agnostic.

\begin{table}[t]
\caption{Engagement-label distribution of selected anchor frames at
$K=5$. ``All'' is the overall DAiSEE distribution for reference.}
\label{tab:anchor-dist}
\centering
\small
\begin{tabular}{lcccc}
\toprule
Source & $L0$ & $L1$ & $L2$ & $L3$\\
\midrule
P-train anchors (109 subjects $\times 5$) & 3 & 25 & 303 & 207\\
P-zero anchors (109 subjects $\times 5$) & 3 & 40 & 263 & 232\\
All DAiSEE clips (n=8571) & 61 & 440 & 4312 & 3758\\
\bottomrule
\end{tabular}
\end{table}

\subsection{Anchor embedding and modulation}
\label{sec:modulation}

We compute the per-subject anchor embedding in the frozen
SigLIP-L feature space as the mean over $A_s$:
\begin{equation}
\mathbf{a}_s \;=\; \frac{1}{|A_s|} \sum_{x \in A_s} f(x), \qquad
f(x) \in \mathbb{R}^{1024}.
\end{equation}

We evaluate three modulation operators applied to each clip $x$
from subject $s$:
\begin{description}\itemsep0pt
\item[\textsc{Sub}] $\phi_{\mathrm{sub}}(x) = f(x) - \mathbf{a}_s$
(1024-d).
\item[\textsc{Cat}] $\phi_{\mathrm{cat}}(x) = [f(x);\; f(x) - \mathbf{a}_s]$
(2048-d): preserves the raw signal alongside the residual,
allowing a linear classifier to use whichever is informative.
\item[\textsc{FiLM}] $\phi_{\mathrm{film}}(x) = (1 + \gamma(\mathbf{a}_s))
\odot f(x) + \beta(\mathbf{a}_s)$, with $\gamma, \beta$ learned linear
maps from the anchor to per-dimension gain and bias. FiLM allows the
model to ignore the anchor where it harms engagement decoding, at
the cost of a small parameter overhead.
\end{description}

We probe each $\phi$ with a class-balanced LR + val-tuned ordinal
thresholds, the strongest recipe identified in
Sec.~\ref{sec:recipe}. The FiLM variant is implemented as a
two-layer MLP head trained with class-balanced cross-entropy
(see~Appendix~\ref{sec:training}).

\subsection{Why we expected \textsc{Dream} to break the ceiling}

The hypothesis was that $f(x) - \mathbf{a}_s$ would represent
``deviation from the subject's neutral-state face,'' which is
intuitively where engagement variation should live. Identity
information would be removed (subtracted off via $\mathbf{a}_s$),
while engagement signal would survive because $\mathbf{a}_s$ is
computed only from neutral, engagement-agnostic frames. The
\textsc{Cat} variant adds the safety net of retaining the raw
signal so that a linear classifier could choose to ignore the
modulation if it hurts. The \textsc{FiLM} variant goes one step
further by letting an MLP head learn how to combine the anchor and
the clip non-linearly.

The reality is more sobering: all three variants fail. We discuss
the result and its implications next.
